% =========================================================================
% A Complete Front-End Functional Verification Environment
% for a SIMT GPGPU: Architecture, Non-Vacuity, and Coverage Closure
%
% Companion to vortex_uvm_paper.tex. That paper argues a METHOD
% (per-instruction SIMT lockstep). This one reports the ENVIRONMENT as an
% engineering artifact, scoped explicitly to front-end functional
% verification, and states its own maturity honestly.
%
% SCOPE DISCIPLINE (do not weaken when editing):
%   "Complete" qualifies the ENVIRONMENT, never the VERIFICATION.
%   The environment is built and self-checking; the campaign has open
%   axes (stimulus volume, error/exception, formal, X-prop) which are
%   reported in Section on maturity rather than omitted.
%
% Numbers are from the banked UCDBs of 2026-08-16 (50 runs per config,
% 0 failures) and the observations register docs/RTL_OBSERVATIONS.md.
%   pdflatex vortex_uvm_frontend_paper.tex  (twice, for refs)
% =========================================================================
\documentclass[conference]{IEEEtran}
\usepackage[utf8]{inputenc}
\usepackage[T1]{fontenc}
\usepackage{amsmath}
\usepackage{booktabs}
\usepackage{array}
\usepackage{url}
\usepackage{listings}
\usepackage{xcolor}
\usepackage{balance}

\lstset{
  basicstyle=\ttfamily\scriptsize,
  breaklines=true,
  frame=single,
  framerule=0.3pt,
  aboveskip=4pt, belowskip=4pt
}

\newcommand{\code}[1]{\texttt{\small #1}}

\begin{document}

\title{A Complete Front-End Functional Verification\\
Environment for a SIMT GPGPU}

\author{\IEEEauthorblockN{Samuel Moussa}
\IEEEauthorblockA{Graduation Project\\
Email: kemetsemiconductors@gmail.com}}

\maketitle

\begin{abstract}
We report a front-end functional verification environment for the Vortex
RISC-V SIMT GPGPU, built and closed by a single engineer. The environment
provides two independent checking layers---byte-exact end-state
equivalence against a golden reference model and per-instruction lockstep
in the writeback domain---and, unusually, \emph{proves both layers can
fail} through committed fault-injection tests that must report an error on
every regression. Coverage reaches 94.7\% total and 98.1\% of covergroup
bins at the primary configuration and 94.6\%/95.8\% at two clusters, with
50 of 50 runs passing at both. Every structural coverage exclusion is
generated per-configuration from register-transfer parameters with a cited
justification, and a blocking merge-time invariant rejects any exclusion
that changes a \emph{covered} bin count; that check found two waivers
which had been silently deleting real coverage. The campaign found real
design defects, including an instruction-set specification deviation and
an IEEE-754 accuracy error. We state the environment's maturity by axis
rather than claiming closure: checker strength and coverage integrity are
strong, whereas stimulus volume, error and exception handling, and formal
methods remain open. We argue that reporting this boundary explicitly is
what separates a verification result from a coverage number.
\end{abstract}

\begin{IEEEkeywords}
functional verification, UVM, SIMT, GPGPU, RISC-V, coverage closure,
non-vacuity, golden reference model
\end{IEEEkeywords}

% =========================================================================
\section{Introduction}

A verification environment is easy to mistake for a verification result.
A bench that compiles, runs a regression, and reports high coverage may
still be checking nothing: if the checker cannot fail, its silence is not
evidence. This paper reports an environment built to make that failure
mode detectable, and states plainly where the resulting campaign is strong
and where it is not.

The design under verification is Vortex, an open-source RISC-V SIMT GPGPU.
It is a demanding target for a small team: a full graphics-processor
hierarchy of clusters, sockets, cores, warps and threads, with a
configurable cache hierarchy, custom SIMT control instructions outside the
base instruction set, and a memory model that is explicitly weakly
coherent.

\subsection{Contributions}
\begin{enumerate}
\item A complete Universal Verification Methodology environment for a SIMT
GPGPU, with role-inverted agents (the design is the bus master; the
environment supplies responders), a golden model integrated through the
direct programming interface, and a passive observability layer bound into
the register-transfer hierarchy so that it scales to any topology without
enumerating paths.
\item \textbf{Non-vacuity as a standing gate.} Both checking layers have
committed negative tests that inject a fault and \emph{must} report an
error. They are re-run after every change, so a refactor that silently
disables checking fails the regression rather than improving it.
\item \textbf{Coverage-exclusion integrity.} Exclusions are generated
per-configuration from elaborated design parameters, each with a cited
justification, and a merge-time invariant blocks any exclusion that
changes a covered-bin count. We report two defects this found, one of
which had been active for weeks.
\item \textbf{An honest maturity assessment} by axis, including the axes
where the campaign is weak, and a categorized roadmap separating remaining
front-end work from back-end sign-off work.
\end{enumerate}

\subsection{What this paper does not claim}
\label{sec:notclaim}
We do not claim the design is verified, nor that it has been stressed.
Neither is supportable. The randomized generator is seed-controlled and
reproducible, but the committed regression runs one seed per profile,
which is directed testing in a random costume. The bus responder always
returns a successful response, so no error-recovery path is exercised.
There is no architectural exception or interrupt stimulus, no formal
property verification, and no X-propagation or randomized-reset analysis.
These are stated here, in the introduction, because a limitation disclosed
only in a final section invites the reader to discover it as a flaw rather
than read it as scope.

% =========================================================================
\section{Environment Architecture}

\subsection{Role inversion}
In a conventional bench the environment drives transactions into a
subordinate design. Here the relationship is inverted: the design is a
program-executing bus master, so the environment's memory agents are
\emph{responders}. Stimulus is therefore selected primarily by choosing
which program to execute, not by randomizing bus transactions. This
distinction matters for how randomization is claimed: the random axis is
the generated program, not the transaction stream.

Four agents are active---host control, device-configuration registers,
and two memory responders (bus-protocol and native)---alongside a passive
status agent and a set of register-transfer probes.

\subsection{Golden model integration}
A functional emulator that ships with the design serves as the reference,
integrated through the direct programming interface. It steps
instruction-by-instruction, models full SIMT architectural state, and
returns retirement records suitable for lockstep comparison.

Two consequences are load-bearing and are reported rather than hidden.
First, the reference is \emph{not independent}: it was co-designed with
the design under verification, so a shared misunderstanding is invisible
to every check. Second, the reference and the design execute the
\emph{same binary}, so a defect in the stimulus itself is common-mode and
cancels---a green run with zero mismatches is compatible with the program
having verified nothing. We return to this in
Section~\ref{sec:silentclass}.

\subsection{Passive observability by binding}
Probes are bound into the register-transfer hierarchy rather than
connected by hierarchical path. A probe bound to a module definition is
instantiated wherever that module is, so an environment written for one
core becomes correct for sixteen with no edit. Probes are strictly
read-only and never gate a verdict; they exist to make coverage
attributable, not to check.

One case required more care. The device-configuration bus is write-only,
so no front-door read exists and ``did this configuration write reach the
register?'' had never been checked. A bound peek-only probe supplies the
read side, enabling a register model with a real read-back check. The
probe never writes: a backdoor write would desynchronize the golden model,
which is fed from the bus monitor.

% =========================================================================
\section{Checking, and Proving the Checkers}
\label{sec:nonvacuity}

\subsection{Two layers}
The primary checker compares final memory state byte-exactly against the
reference model. It is bidirectional: a forward pass compares every
location the design wrote, and a reverse pass compares result-region
locations the design \emph{failed} to write. The reverse pass exists
because a dropped store is invisible to the forward pass---there is no
transaction to compare.

The secondary checker performs per-instruction lockstep in the writeback
domain, comparing program counter, destination register and value at
retirement.

\subsection{Non-vacuity as a standing gate}
Both layers have committed negative tests. One injects a wrong value into
a result location; the other drops a store. Both must report an error, and
both are re-run after every change to the environment. When the scoreboard
was later rebuilt onto a single data source, the negative tests were
confirmed still failing at the same address---which is what makes the
refactor trustworthy rather than merely green.

We regard this as the most important property of the environment. A
checker that has never been observed to fail is indistinguishable from an
absent one, and most published verification results do not report
evidence that their checkers can fail at all.

\subsection{The silent-failure class}
\label{sec:silentclass}
Because both models execute the same binary, a stimulus defect cancels.
This is not hypothetical. Two instances were found in this project.

In the first, a compute kernel derived a tile geometry from a compile-time
template parameter rather than from the runtime device configuration. The
kernel was scaled to a larger topology, the geometry silently did not
follow, and both models computed the same wrong answer and agreed. The
end-state compare passed and coverage filled.

In the second, found while preparing the cache-hierarchy experiment
reported in Section~\ref{sec:results}, a kernel's phase-enable flags were
discarded by the build system: a command-line variable overrode the
makefile assignment that added them. The kernel compiled with every phase
disabled, executed, compared byte-exactly against the reference, and
passed---having exercised nothing. The defect was invisible in the
simulation log and visible only by comparing the two compiler invocations
directly.

The general lesson is that \emph{only what differs between the two models
is validated}; everything shared requires an independent assertion. We now
treat the build path as part of the verified surface.

% =========================================================================
\section{Coverage Methodology}

\subsection{Exclusions as generated artifacts}
Structural coverage exclusions are not maintained by hand. A generator
emits them per-configuration from elaborated design parameters, so a
waiver that is legitimate at one topology is not applied at another where
the same logic is reachable. Each carries a citation to the
register-transfer source that makes it structurally dead.

Two examples illustrate the discipline. The instruction cache is
instantiated with its write port disabled, so its entire write-data path
is elaborated but undrivable; this accounts for 26\% of the design's
toggle gap from a single subtree, and is waived with the instantiation
cited. A positive control confirms the diagnosis: the corresponding
response-data path toggles fully on all bits. By contrast, the upper bits
of a 44-bit retirement identifier are \emph{not} waived, because reaching
them requires $2^{18}$ or more instructions---infeasible is not the same
as structurally dead, and only the latter justifies removal from the
denominator.

\subsection{A blocking invariant on the merge}
Exclusions are applied in two stages separated by justification class.
Third-party intellectual property is excluded because it is not the design
under verification, and legitimately removes covered bins. Structural
exclusions must remove \emph{only} unreachable bins. The merge therefore
asserts a hits-invariant: after structural exclusions, no covered-bin count
may change. Violation fails the merge.

This check found two real defects on its first execution. One was a waiver
written the same day. The other had been silently discarding a covered
condition term for weeks. The root cause in both cases is a tool property
worth recording: focused expression coverage cannot waive a single input
term of a condition---waiving one term discards the whole expression,
including its covered rows.

\subsection{What the total metric means}
The tool's total coverage is the unweighted mean of seven categories,
each contributing equally regardless of bin count. The lowest category is
therefore the largest lever, independent of its size. We report the
per-category breakdown rather than the total alone, because the total can
be moved by a category with five bins as easily as by one with 400,000.

% =========================================================================
\section{Results}
\label{sec:results}

\begin{table}[t]
\caption{Banked coverage (per configuration; never blended)}
\label{tab:fe-coverage}
\centering
\small
\begin{tabular}{lrr}
\toprule
Metric & 1CL/1C/4W/4T & 2CL/2C/4W/4T \\
\midrule
Covergroup bins (raw) & 370/377 = 98.1\% & 989/1032 = 95.8\% \\
Covergroup (weighted) & 99.8\% & 99.5\% \\
Statement & 98.1\% & 98.3\% \\
Branch & 95.1\% & 95.7\% \\
Condition & 90.4\% & 88.8\% \\
Toggle & 82.8\% & 80.5\% \\
Assertion & 96.9\% & 98.9\% \\
Directive & 100.0\% & 100.0\% \\
\midrule
\textbf{Total} & \textbf{94.7\%} & \textbf{94.6\%} \\
Runs passing & 50/50 & 50/50 \\
Coverage instances & 2\,256 & 8\,275 \\
\bottomrule
\end{tabular}
\end{table}

Table~\ref{tab:fe-coverage} reports the two banked configurations. Banks
are never blended across configurations: instance counts and signal widths
differ, so a merged database is not meaningful.

\subsection{Directed stimulus for structural targets}
Where a coverage gap resisted the existing suite, kernels were written
against the specific missing terms and verified against those terms
\emph{before} entering the regression. One produced genuine internal
back-pressure and covered thirteen of twenty-four missing condition terms
at the primary configuration; another created a second concurrent memory
stream by executing a large resident instruction footprint interleaved with
data traffic.

\subsection{A negative result on memory pressure}
An intuition worth publishing because it is wrong: increasing memory
pressure does not increase contention coverage. Two independent
experiments enlarged a kernel's working set by a factor of sixteen. One
went from thirteen covered condition terms to eleven, losing a target
outright; the other went from four to zero.

The mechanism is that the terms in question require two requests live at
one port in the same cycle. A working set that exceeds cache capacity
sends nearly every access to main memory, every thread blocks, and
issue stops. Requests still arrive, but spread thin---each drained long
before the next. Total traffic rises while \emph{concurrent} traffic
falls. Contention is produced by issue density, not by miss volume. Both
kernels carry this reasoning in their source, because the
well-intentioned edit that would undo it looks like an optimization.

\subsection{Exercising the shared cache hierarchy}
The optional second- and third-level caches are instantiated in bypass
form unless explicitly enabled, and no regression had ever enabled them:
the suite had no path from its configuration to the compile. After
supplying that path, a reuse-based kernel---touch a working set sized to a
level, then re-read it---covered the hit path of all four shared-cache
instances for the first time, with 13\,000 to 15\,000 hits each, while
comparing 196\,868 memory words byte-exactly against the reference with
zero errors. This is the largest end-state comparison in the suite by a
factor of six.

\subsection{Defects found}
The campaign found defects in both the design and the reference model,
maintained in an evidence-cited register with explicit dispositions. The
design-side findings include a deviation from the instruction-set
specification in indirect-jump target computation and an IEEE-754 accuracy
error of one unit in the last place in hardware square root. Several
findings are observability limits or micro-architectural quirks rather
than defects, and are recorded as such; a register that only ever records
bugs is not being maintained honestly.

% =========================================================================
\section{Maturity by Axis}
\label{sec:maturity}

We assess the campaign by axis rather than reporting a single readiness
figure, because the axes are not substitutable: no amount of checker
strength compensates for absent error-path stimulus.

\begin{table}[t]
\caption{Front-end verification maturity, self-assessed}
\label{tab:maturity}
\centering
\small
\begin{tabular}{ll}
\toprule
Axis & Status \\
\midrule
Environment architecture & Strong \\
Checker strength / non-vacuity & Strong; above typical practice \\
Coverage integrity \& waiver audit & Strong \\
Defect-discovery record & Strong \\
Reference-model quality & Good, but \emph{not independent} \\
Directed stimulus & Good \\
Configuration breadth & Weak (three points) \\
Coverage-model provenance & Moderate; not specification-traced \\
Random stimulus volume & \textbf{Weak} (one seed per profile) \\
Error / exception verification & \textbf{Absent} \\
Formal property verification & \textbf{Absent} \\
X-propagation / reset randomization & \textbf{Absent} \\
\bottomrule
\end{tabular}
\end{table}

Table~\ref{tab:maturity} is deliberately unflattering in its lower half.
The environment---the part that took the engineering---is substantially
complete. The campaign that runs on it is not, and the deficit is
concentrated in stimulus rather than infrastructure. That is an
encouraging shape, because stimulus volume is bought with machine time
once the generator is seed-controlled, which it now is.

A second structural limit deserves emphasis. The reference model is not
independent of the design. An independent scalar cross-check was performed
against a third model and agreed on all 11\,076 base-instruction
retirements, but that model has no SIMT semantics and cannot execute the
design's custom control instructions. The SIMT axis therefore has no
independent reference and will not get one from that direction. This is a
ceiling, not an unfinished task.

% =========================================================================
\section{Toward Tape-Out}
\label{sec:roadmap}

The two lists below are separated because they support different claims.
Front-end items strengthen the claims in this paper; back-end items are
prerequisites for a different claim entirely---that the design is
manufacturable and will work in silicon. Nothing here speaks to the
latter.

\subsection{Remaining front-end work}
\begin{enumerate}
\item \textbf{Stimulus volume.} Scale the seeded generator from one seed
per profile to hundreds. Highest return in this list; costs machine time,
not engineering.
\item \textbf{Error injection} on the bus responder (error and decode
responses), opening the untested error-recovery axis.
\item \textbf{Exception and interrupt stimulus}, designed around the
observation that this design terminates by deasserting a busy signal
rather than by executing a breakpoint---the assumption that blocked
earlier attempts.
\item \textbf{Formal property verification} on arbitration-heavy control
where simulation is weakest and the state space is small.
\item \textbf{X-propagation and randomized reset}, to reach
uninitialized-state behavior that a two-state flow cannot see.
\item \textbf{Configuration-matrix breadth} beyond three sampled points.
\item \textbf{Specification-traced coverage model}, so closure measures
the specification rather than the model.
\end{enumerate}

\subsection{Back-end and ASIC sign-off work}
\begin{enumerate}
\item Gate-level simulation, then with back-annotated timing.
\item Static timing analysis across process, voltage and temperature
corners.
\item Design-for-test: scan, automatic test-pattern generation and
fault-coverage sign-off; memory self-test.
\item Clock- and reset-domain-crossing analysis, structural and
functional.
\item Lint and structural sign-off against a synthesis rule set.
\item Low-power verification: power intent, retention and isolation.
\item Physical closure: place and route, extraction, signal and power
integrity.
\item Equivalence checking across source, netlist and post-layout netlist.
\item Post-silicon bring-up, characterization, and a path from silicon
failures back into this regression.
\end{enumerate}

% =========================================================================
\section{Conclusion}

We reported a front-end functional verification environment for a SIMT
GPGPU in which both checking layers are proven able to fail, every
structural coverage exclusion is generated from design parameters with a
citation, and a blocking merge-time invariant prevents a waiver from
consuming real coverage. Coverage closed at 94.7\% total and 98.1\% of
covergroup bins at the primary configuration, with all runs passing at
both configurations, and the campaign found real defects in the design.

The environment is substantially complete; the campaign is not, and we
have said where. Our broader argument is that this distinction is the
substance of a verification result. A coverage number is a measurement of
a model, and it is only as meaningful as the evidence that the checkers
could have failed, that the waivers removed nothing real, and that the
stimulus exercised what it claims. Where we could not supply that
evidence, we have marked the axis open rather than closing it by
assertion.

\balance
\end{document}
